
% gloss: acquisition timing
\section{Acquisition timing and order properties of the core}\label{sec:timing}

Profile steps can add held tests and never drop them
(Definition~\ref{def:profile-steps}). For sequences with profile steps,
acquisition timing can change fine reachability.
Section~\ref{sec:acquisition} gives a polynomial-time signed-core test in
the single-word class and a projected-component criterion for the full
class.
The hardness reductions of Section~\ref{sec:hard} for the Markov model use
a fixed profile, so their instances involve no acquisition timing.
Section~\ref{sec:core-order} shows when the profile order alone decides
reachability.

% ---------------------------------------------------------------------------
\subsection{Acquisition timing}\label{sec:acquisition}

An \defterm{acquisition} is a profile step whose cover adds one held test,
keeping the table support, partitions and parameter vector.
\defterm{Acquisition timing} is whether a test is acquired before or after
a numerical revision. The sign example of Section~\ref{sec:directed}
shows this effect.
Both endpoint parameter vectors pass $\mathsf{yy}\sim\mathsf{zz}$, yet $B$
is finely reachable from~$A$ and not from~$A^+$: $A$ can take $y$
through~$0$ before acquiring the test, whereas $A^+$ holds the test from
the start.

\begin{theorem}[Acq{}uisition timing, single-word class]\label{thm:acquisition}
Fix a shared-probe model in the single-word class, and feasible
configurations $A$, $A^+$ and~$B$, where $A^+$ has the parameter vector
of~$A$ and $p_A\preceq p_{A^+}\preceq p_B$. Then $B$ is finely reachable
from $A$ and not from $A^+$ if and only if
$\SC_{p_A}(\alpha_A)=\SC_{p_A}(\alpha_B)$ and
$\SC_{p_{A^+}}(\alpha_A)\neq\SC_{p_{A^+}}(\alpha_B)$. With explicit binary
inputs promised feasible, two signed-core comparisons, each computing both
endpoint cores by Theorem~\ref{thm:certificates}(a), decide this in
polynomial time.
\end{theorem}

Under the hypotheses of Theorem~\ref{thm:acquisition}, the timing effect
occurs exactly when the endpoint signed cores agree under the source tests
and differ under the tests held after acquisition. The acquired test may be a value test. An acquisition never enlarges the set of reachable targets
(Corollary~\ref{cor:refine}).

A held word-against-pair test can couple label probabilities and
eigenvalues (Lemma~\ref{lem:compile}(b)), and Definition~\ref{def:system}(e)
defines $\Phi_p$ for the single-word class. For the full class, we instead
use the projected tables $\bar P_{A,m}=\theta_{A,m}Q_{A,m}$ $(m\in M)$
from Definition~\ref{def:profile-steps}. For an attained profile~$p$, the
\defterm{projected rational closure} $\mathcal K_p\subseteq[0,1]^{M\times X\times X}$ is the sup-norm
closure of the projected tables of feasible rational configurations at~$p$.
The projected rational closure is compact. It can be strictly smaller than
the locus of projected tables of feasible real configurations at the same
profile, even with one probe (Proposition~\ref{prop:one-probe}).

\begin{theorem}[Projected components]\label{thm:projected-components}
Fix a shared-probe model, with word tests of both shapes allowed, and
feasible configurations $A$ and~$B$.
\begin{enumerate}
\item[(a)] If $p_A\preceq p_B$, then $\theta_B$, $\alpha_B$ and the
register tables of $A$ form a feasible configuration $A_B$ at~$p_A$, and
$B$ is finely reachable from $A$ if and only if $A_B$ is.
\item[(b)] $B$ is finely reachable from $A$ if and only if
$p_A\preceq p_B$ and $\bar P_A,\bar P_B$ lie in one connected component
of~$\mathcal K_{p_A}$.
\item[(c)] If $p\preceq q$ are attained, then
$\mathcal K_q\subseteq\mathcal K_p$.
\item[(d)] For a feasible $A^+$ with
$p_A\preceq p_{A^+}\preceq p_B$, $B$ is finely reachable from $A$ and
not from $A^+$ if and only if $\bar P_A,\bar P_B$ share a component of
$\mathcal K_{p_A}$ and $\bar P_{A^+},\bar P_B$ do not share one
of~$\mathcal K_{p_{A^+}}$.
\end{enumerate}
\end{theorem}

As in Theorem~\ref{thm:directed}, first check $p_A\preceq p_B$: if this
order fails, no sequence exists at any tolerance. Otherwise, the verdict
is whether the two projected tables lie in the same connected component
of~$\mathcal K_{p_A}$ (Lemma~\ref{lem:compact-chain};
Appendix~\ref{app:timing-proofs}).
In~(d), take the directed interval $p_A\preceq p_{A^+}\preceq p_B$ and a
feasible acquisition that keeps the parameter vector fixed. Then the
timing effect occurs exactly when the endpoints' projected tables share a
component before acquisition and lie in different components afterward.
In the sign example, $\mathcal K_{p_A}$ is a segment of tables indexed
by~$y$, and $\mathcal K_{p_{A^+}}$ is the two tables with $y=\pm1/4$. A
feasible acquisition need not split a component.

We give no algorithm for the components of~$\mathcal K_p$. Because
$\mathcal K_p$ is a closure of rational projected tables in the real
projected-table space, deciding its components is related to Mazur's
question on the real topology of rational points~\cite{mazur}.

\begin{proposition}[Source-profile scheduling]\label{prop:source-schedule}
Fix a shared-probe model, with word tests of both shapes allowed, feasible
configurations $A$ and $B$ with $B$ finely reachable from~$A$,
and attained profiles $p_A\preceq q_1\preceq\dots\preceq q_k\preceq p_B$.
For every rational $\varepsilon>0$, a reconfiguration sequence at
tolerance $\varepsilon$ leads from $A$ to the configuration $A_B$ of
Theorem~\ref{thm:projected-components}(a), then through configurations at
$q_1,\dots,q_k$ in order, each with label probabilities $\theta_B$ and
eigenvalues $\alpha_B$, to~$B$.
\end{proposition}

Thus the label-probability and eigenvalue revisions can be completed at
the source profile, before any test is acquired (Theorem~\ref{thm:projected-components}(a)).
Register rewrites may occur in the later phase. The sequence can pass
through the target's full parameter vector at a prescribed intermediate
profile whose support equals the target support.

Barrett and Weld's account of Korf's subgoal taxonomy~\cite{korf}
distinguishes goals that can be achieved and protected in a suitable order
from goals for which no such order
works~\cite[Sections~3.2--3.3]{barrettweld}. In Hauser and Latombe's
multimodal planning, a transition point in the intersection of two modes
can lie in a component inaccessible from the start in the old mode or from
the goal in the new mode~\cite[Sections~2.3--2.5]{hauserlatombe}. Hu, Iyer
and Christensen connect within-manifold roadmaps at
intersections~\cite[Section~IV.B]{huiyerchristensen}. To our knowledge,
the before-and-after criterion for projected rational closures in this
model has not been stated.

% ---------------------------------------------------------------------------
\subsection{Order properties of the core}\label{sec:core-order}

% gloss: G17
\begin{theorem}[No nonzero calibration]\label{thm:no-calibration}
Fix a shared-probe model in the single-word class and a profile~$p$. For
every rational $\alpha$ satisfying the equations of~$\Phi_p$,
$\SC_p(\alpha)=0$ if and only if $p$ holds no calibration
$\alpha_{mi}=c$ with $c\neq0$. Hence, if $A$ is feasible and $p_A$ holds
no such calibration, a feasible $B$ is finely reachable from $A$ if and
only if $p_A\preceq p_B$.
At a fixed profile with no nonzero calibration, any two feasible
configurations are mutually finely reachable. Setting all eigenvalues of
$A$ to zero while keeping its label probabilities and register tables gives
a feasible configuration at the same profile, mutually finely reachable
with~$A$.
\end{theorem}

For a vector satisfying the compiled equations, zero calibrations are
inactive on its support. If no nonzero calibration is held, the active
equations have equal degrees on their two sides, so unit weights peel the
entire support (Proposition~\ref{prop:empty-core}). A held nonzero
calibration keeps its coordinate in the core. Without
$z=1/4$, the sign example's test gives $2\eta_y=2\eta_z$, met by
$\eta=(1,1)$.
The empty-core conclusion in the absence of held nonzero calibrations can
fail for arbitrary binomial tests: under
$(xy)^2=xy$ and $|x|+|y|\le2$, the cores of the endpoints $(1,1)$ and
$(-1,-1)$ both have support $\{x,y\}$, with signed cores $(+,+)$ and
$(-,-)$. Moreover, $\alpha$ must satisfy~$\Phi_p$: with $z=0$ held and
the word test unheld, the formal support $\{z\}$ is irreducible, but no
eigenvalue vector passing the calibration has that support.

% gloss: G18
With budgets and metric fixed, adding tests only splits the classes of the
points that satisfy them, and a new test can also exclude points
(Corollary~\ref{cor:refine}).

% gloss: G19
By Lemmas~\ref{lem:restriction} and~\ref{lem:irreducible}, the core is
an interior operator on supports, and its removed-coordinate map
$Z\mapsto J\setminus\core(J\setminus Z)$ is a closure operator whose
closed sets are the complements of irreducible supports
(Proposition~\ref{prop:interior}). The core can strictly shrink an
attained support: at the sign example's source profile, the attained
support $\{y,z\}$ has core~$\{z\}$.
